\documentclass[10pt,a4paper]{amsart}
\usepackage[margin=19mm]{geometry}
\usepackage{amsmath,amssymb,mathtools}
\usepackage[scr=rsfs,cal=euler]{mathalfa}
\usepackage{xfrac}
\usepackage[unicode,hidelinks,bookmarks=false]{hyperref}
\newcommand{\ie}{i.e.\ }
\newcommand{\cf}{cf.\ }
\newcommand{\resp}{respectively\ }
\newcommand{\Subsection}{\subsection}
\providecommand{\nobreakdash}{\nobreak\hbox{-}}
\setlength{\emergencystretch}{2em}
\multlinegap0pt
\usepackage{bm}
\usepackage{bbm}
\usepackage{dsfont}
\usepackage{xspace}
\usepackage{cancel}
\usepackage{mathtools}
\usepackage{amsthm}
\makeatletter
\@ifundefined{theorem}{\newtheorem{theorem}{Theorem}}{}
\@ifundefined{claim}{\newtheorem{claim}[theorem]{Claim}}{}
\makeatother
\newcommand{\eps}{\varepsilon}
\renewcommand{\ge}{\geqslant}
\renewcommand{\le}{\leqslant}
\title[An edge-spectral supersaturation of Mubayi's theorem for col]{An edge-spectral supersaturation of Mubayi's theorem for color-critical graphs}
\author{Hongzhang Chen, Yongtao Li}
\date{Source: arXiv:2607.01073v2 (CC BY 4.0)}
\begin{document}
\maketitle

\noindent\emph{Source and licence.} An edge-spectral supersaturation of Mubayi's theorem for color-critical graphs. Version arXiv:2607.01073v2 (CC BY 4.0), distributed under the Creative Commons Attribution 4.0 International licence, as recorded in the arXiv licence metadata for this exact version and shown on the abstract page https://arxiv.org/abs/2607.01073v2. Full rights notice: licenses/arxiv-2607.01073-CC-BY-4.0.txt.\par\smallskip

\noindent\textit{Editorial scope.} Counted unit: the Claim whose source label is \texttt{clm:perron-almost-constant} in the article (source0.tex lines 1329--1332, source label \texttt{clm:perron-almost-constant}), delivered together with the article's own complete original proof of it (source0.tex lines 1334--1410, one \texttt{proof} environment of 77 lines). No mathematics is written, repaired or re-derived: statement and proof are the article's own text line for line. Required context retained uncounted: clm:n-size (lines 952--955); clm:balanced-partition (lines 994--1003). Every definition, display and label of those lines is retained unchanged. Omitted from this package and not redistributed: 1--951, 956--993, 1004--1328, 1333--1333, 1411--2007. Nothing inside a retained excerpt is omitted or altered: every retained body is delivered exactly as the article's lines are written, so no omission receipt of any kind accompanies this package. Citations: the retained excerpts carry no citation command at all, so no bibliography entry of the article is transcribed and no reference is redistributed. Reference adaptation: none was needed and none was made; every reference inside the delivered unit resolves inside it, so no unresolved reference or citation is produced. Printed slips kept literal: no printed word, symbol, hypothesis or number of the counted statement or of its proof was altered. Layout and class: the article's own class is \texttt{article} with packages including geometry, amsmath, amssymb, amsthm, mathtools, bm, enumitem, microtype, hyperref, tikz; the wrapper is the lane's pinned \texttt{amsart} editorial wrapper at 10pt on A4, which restates the theorem-like environments the retained text needs and adds the article's own macro definitions that the retained text needs (\texttt{\textbackslash eps}, \texttt{\textbackslash ge}, \texttt{\textbackslash le}), with extra wrapper packages (bm, bbm, dsfont, xspace, cancel, mathtools). The delivered numbering is the unit's own. Assets: no retained span contains a figure environment, a graphics-inclusion command, a PDF-page inclusion, a table or a TikZ picture; no image file is copied into this package, the package declares no asset dependency and no image file is redistributed with it. Licence: version arXiv:2607.01073v2 of this article is distributed under the Creative Commons Attribution 4.0 International licence, as recorded in the arXiv licence metadata for this exact version and shown on the abstract page https://arxiv.org/abs/2607.01073v2. A full rights notice is delivered at licenses/arxiv-2607.01073-CC-BY-4.0.txt. Characters: every retained excerpt is pure ASCII, so the delivered engine, XeLaTeX with its default Latin Modern text font, reports no missing character.
\begin{claim}\label{clm:n-size}
Let $ \psi=\sqrt{\frac{2h}{1-1/r}}$. Then 
$ (1-\eps_1)\psi\le n\le (1+\eps_1)\psi.$
\end{claim}

\begin{claim}\label{clm:balanced-partition}
We have $  e(H)\ge \frac{r-1}{2r}n^2-3\eps_1^2n^2$, and 
the partition satisfies
\[
  \sum_{i=1}^r e(H[V_i])\le \eps_1^2n^2, \quad \text{and~} \quad 
    \left||V_i|-\frac nr\right|\le 3\eps_1n
  \quad
  \text{for every }i\in[r]. 
\]
\end{claim}

\begin{claim}\label{clm:perron-almost-constant}
Every $u\in\bigcup_{i=1}^r V_i^{(1)}$ satisfies
$x_u\ge(1-16f^2\eps_1)x_{u^\ast}$.
\end{claim}

\begin{proof}
Throughout we use $|V_j|\le(\tfrac1r+3\eps_1)n$
(Claim~\ref{clm:balanced-partition}). By Claim~\ref{clm:n-size},
$\lambda(H)\ge\sqrt{2(1-\tfrac1r)h}=(1-\tfrac1r)\psi$ with
$\psi\ge(1-\eps_1)n$, so
$\lambda(H)\ge(1-\tfrac1r-\eps_1)n$; since $r\ge3$, this
gives $\lambda(H)>\tfrac n2$.

\emph{Non-neighbors.} Fix
$u\in V_i^{(1)}=V_i\setminus(W_i\cup S^{(1)})$, so
$d_{V_i}(u)<4\eps_1n$ and $d_H(u)>(1-\tfrac1r-4\eps_1)n$.
For $j\ne i$, splitting off the parts $V_i$ and $V_s$
($s\ne i,j$) gives
$d_{V_j}(u)=d_H(u)-d_{V_i}(u)-\sum_{s\ne i,j}d_{V_s}(u)
>(\tfrac1r-3f\eps_1)n$ (using $f\ge r+1$). 
Hence
$
  \sum_{j\ne i}\bigl(|V_j|-d_{V_j}(u)\bigr)
  <3(f+1)(r-1)\eps_1n<4(f^2-1)\eps_1n$.

\emph{Same part.} Let $u_1,u_2\in V_i^{(1)}$ with
$x_{u_1}\ge x_{u_2}$. From the eigenvalue equation, we have 
$$\lambda(H)(x_{u_1}-x_{u_2})
=\sum_{v\in V(H)}(a_{u_1v}-a_{u_2v})x_v
\le\sum_{v\in N_H(u_1)\setminus N_H(u_2)}x_v.$$
Each such $v$ lies either in $N_{V_i}(u_1)$ (at most
$d_{V_i}(u_1)<4\eps_1n$ vertices) or, for some $j\ne i$, in
$V_j\setminus N_{V_j}(u_2)$, and the latter total is at
most $4(f^2-1)\eps_1n$ by the non-neighbor bound applied to
$u_2$. As each $x_v\le x_{u^\ast}$, we get
$\lambda(H)(x_{u_1}-x_{u_2})\le4f^2\eps_1n\,x_{u^\ast}$, and
$\lambda(H)>\tfrac n2$ then yields
$x_{u_1}-x_{u_2}\le8f^2\eps_1x_{u^\ast}$. Thus, the Perron
coordinates inside one $V_i^{(1)}$ differ by at most
$8f^2\eps_1x_{u^\ast}$.

\emph{Location of $u^\ast$.} Since
$\lambda(H)x_{u^\ast}=\sum_{v\in N_H(u^\ast)}x_v
\le d_H(u^\ast)x_{u^\ast}$, we have
$d_H(u^\ast)\ge\lambda(H)\ge(1-\tfrac1r-\eps_1)n$, so
$u^\ast\notin S^{(1)}$. As $\bigcup_iW_i\subseteq S^{(1)}$,
this places $u^\ast\in\bigcup_iV_i^{(1)}$, say
$u^\ast\in V_1^{(1)}$. 
Applying the same-part bound established above,  
we then obtain 
$x_u\ge(1-8f^2\eps_1)x_{u^\ast}$ for all $u\in V_1^{(1)}$.

\emph{Other parts.} Suppose some
$u_0\in\bigcup_{i\ge2}V_i^{(1)}$, say $u_0\in V_2^{(1)}$,
had $x_{u_0}<(1-16f^2\eps_1)x_{u^\ast}$. Since coordinates
in $V_2^{(1)}$ differ by at most $8f^2\eps_1x_{u^\ast}$,
every $u\in V_2^{(1)}$ then satisfies
$x_u<(1-8f^2\eps_1)x_{u^\ast}\le(1-8r^2\eps_1)x_{u^\ast}$,
the last step using $f\ge r$. We now bound
$\lambda(H)x_{u^\ast}=\sum_i\sum_{v\in N_{V_i}(u^\ast)}x_v$
part by part. For $V_2$, the at most
$|V_2|\le(\tfrac1r+3\eps_1)n$ neighbors in $V_2^{(1)}$ each
contribute at most $(1-8r^2\eps_1)x_{u^\ast}$, while the at most
$|W_2\cup S^{(1)}|\le2\eps_1n$ neighbors outside contribute
at most $ x_{u^\ast}$, so
$\sum_{v\in N_{V_2}(u^\ast)}x_v
\le \left((1-8r^2\eps_1)(\tfrac1r+3\eps_1)n+2\eps_1n \right) x_{u^\ast}$. For $V_1$, since
$u^\ast\notin W_1$ we have $d_{V_1}(u^\ast)<4\eps_1n$, so
this part contributes at most $4\eps_1n\,x_{u^\ast}$; each
$V_i$ with $i\ge3$ contributes at most
$(\tfrac1r+3\eps_1)n\,x_{u^\ast}$. 
Adding these, it follows that 
\[
  \lambda(H)
  \le(r-1-8r^2\eps_1)\bigl(\tfrac1r+3\eps_1\bigr)n+6\eps_1n
  \le\bigl(1-\tfrac1r-4r\eps_1\bigr)n,
\]
where the last step expands the coefficient to
$1-\tfrac1r-(5r-3)\eps_1-24r^2\eps_1^2$ and uses $r\ge3$.
This contradicts $\lambda(H)\ge(1-\tfrac1r-\eps_1)n$, so no
such $u_0$ exists, and the claim follows.
\end{proof} 

\end{document}
